\section{
    Almost simply typed reinterpretable language: STLC-R
}\label{STLC-R}

We start with a simpler language to iron out the details
and actually figure out how to properly reinterpret
terms internally before moving on to dependently typed
system. Our starting point is simply typed lambda
calculus with products, which we call <<STLC>>
throughout this paper. As discussed before, we will add
self-reinterpretation capabilities to it to enable
reinterpretation in any <<model>> of STLC specifiable
inside our language. (TODO: enter cartesian categories
and Curry-Howard-Lambek correspondence) Here's the thing
we \textit{wish} to do with it, where
$t\,@\,\mathcal{M}$ is an operator of reinterpretation
of term $t$ in internal model $\mathcal{M}$:
\begin{multline*}
    (\lambda f x y. (f x, y))\,@\,\langle(\circ),
        \mathtt{pair},\mathtt{proj}_1,\mathtt{proj}_2,
        \mathtt{lam},\mathtt{eval},\ldots\rangle=\\=
    \mathtt{lam}\,(\mathtt{lam}\,(\mathtt{lam}\,(
        \mathtt{pair}\,
            (\mathtt{eval}\,\circ\,\mathtt{pair}\,
                (\mathtt{proj}_2\circ\mathtt{proj}_1
                    \circ\mathtt{proj}_1
                )\,
                (\mathtt{proj}_2\circ\mathtt{proj}_1)
            )\,
            \mathtt{proj}_2
    )))
\end{multline*}

Where $\mathtt{lam}$,\,$\mathtt{pair}$ and
$\mathtt{eval}$ are correspondingly constructor of an
exponential object, constructor of a product and
evaluation morphism, in the same way they are specified
in the definition of cartesian closed categories.

There are a few things to ponder:

\begin{itemize}
    \item How to reinterpret free variables?
    \item How to type universes?
    \item We're in a simply typed setting, so there's no
        way to enforce that universe satisfies coherence
        conditions on their operations. How to structure
        computation in such a way that it doesn't break
        Church-Rosser property?
\end{itemize}

To answer the first question, we actually disallow (most
of) free variables in a terms that are to be
reinterpreted since it is not clear how to lift
arbitrary \textit{value} into current universe, so we
have to enforce some kind of a typing discipline for
introducing terms amenable to reinterpretation. In a
sense, reinterpretation is a strong form of
polymorphism, so let's structure our grammar and typing
rules similarly to how it's done in Hindley-Milner type
systems today and see how far we can go; grammar rules
are collected in a figure \ref{stlc-r-grammar}.

\begin{figure}
\centering
\begin{bnf}
    $S,T,C,P,E,TT$ : \textsf{MonoType} ::=
    | $\alpha$ : type variable
    | $\top$ : unit type
    | $S\times T$ : product type
    | $S\to T$ : arrow type
    ;;
    $\sigma$ : \textsf{PolyType} ::=
    | $T$ : monomorphic type
    | $\forall\alpha.\sigma$ : forall type
    ;;
    $s,t$ : \textsf{Term} ::=
    | $x$ : variable
    | $()$ : singleton
    | $(s, t)$ : pair
    | $\pi_1\,t$ : first projection
    | $\pi_2\,t$ : second projection
    | $\lambda x. t$ : abstraction
    | $t\,s$ : application
    | \texttt{let}\,$x=s$\,\texttt{in}\,$t$
        : polymorphic binding
    | $x^\mathcal{M}$ : reinterpretation
    ;;
    $\mathcal{M}$ : \textsf{Model} ::=
    | $\langle T, t, \ldots \rangle$
        : model def. (see below)
\end{bnf}
\caption{STLC-R grammar}
\Description{}
\label{stlc-r-grammar}
\end{figure}

\begin{figure}
\centering
\begin{mathpar}
    \inferrule{x:\sigma\in\Gamma}{\Gamma\vdash x:\sigma}
    \and
    \inferrule{}{\Gamma\vdash ():\top}
    \and
    \inferrule{\Gamma\vdash s:S\\\Gamma\vdash t:T}
        {\Gamma\vdash (s,t):S\times T}
    \and
    \inferrule{\Gamma\vdash t:S\times T}
        {\Gamma\vdash \pi_1\,t:S}
    \and
    \inferrule{\Gamma\vdash t:S\times T}
        {\Gamma\vdash \pi_2\,t:T}\\
    \and
    \inferrule{\Gamma,x:S\vdash t:T}
        {\Gamma\vdash\lambda x.t:S\to T}
    \and
    \inferrule{\Gamma\vdash t:S\to T\\\Gamma\vdash s:S}
        {\Gamma\vdash t\,s:T}
    \and
    \inferrule{\Gamma\vdash t:\sigma}
        {\Gamma\vdash t:\forall\alpha.\sigma}
    \and
    \inferrule{\Gamma\vdash t:\forall\alpha.\sigma}
        {\Gamma\vdash t:\sigma[\alpha\coloneqq T]}
    \and
    \inferrule{
        \mathrm{fv}(\sigma)=\varnothing\\
        \forall\Gamma\vdash s:\sigma\\
        \forall\Gamma,x:\sigma,\Delta\vdash t:T
    }{\forall\Gamma,\Delta\vdash
        \mathtt{let}\,x=s\,\mathtt{in}\,t:T}
    \and
    \inferrule{
        x:\forall\vec{\alpha}.T\in\forall\Gamma\\
        \forall\Gamma,\Delta\vdash
            \mathcal{M}\,\mathrm{model}
    }{\forall\Gamma,\Delta\vdash x^\mathcal{M}
        :\llbracket T\rrbracket^\mathcal{M}_{
            \left[\vec{\alpha}\coloneqq\vec{S}\right]}}
    \and
    \inferrule{
        \mathrm{fv}(C)\cup\mathrm{fv}(P)\cup
            \mathrm{fv}(E)\subseteq\{0,1\}\\
        \Gamma\vdash\mathtt{id}:
            \forall\alpha.C[\alpha,\alpha]\\
        \Gamma\vdash\mathtt{comp}:
            \forall\alpha\beta\gamma.
            C[\beta,\gamma]\to C[\alpha,\beta]\to
            C[\alpha,\gamma]\\
        \mathrm{fv}(TT)=\varnothing\\
        \Gamma\vdash\mathtt{cnst}:
            \forall\alpha.C[\alpha,TT]\\
        \Gamma\vdash\mathtt{pair}:
            \forall\alpha\beta\gamma.
            C[\gamma,\alpha]\to C[\gamma,\beta]\to
            C[\gamma,P[\alpha,\beta]]\\
        \Gamma\vdash\mathtt{proj}_1:
            \forall\alpha\beta.
            C[P[\alpha,\beta],\alpha]\\
        \Gamma\vdash\mathtt{proj}_2:
            \forall\alpha\beta.
            C[P[\alpha,\beta],\beta]\\
        \Gamma\vdash\mathtt{lam}:
            \forall\alpha\beta\gamma.
            C[P[\gamma,\beta],\alpha]\to
            C[\gamma,E[\alpha,\beta]]\\
        \Gamma\vdash\mathtt{eval}:
            \forall\alpha\beta.
            C[P[E[\alpha,\beta],\beta],\alpha]
    }{\Gamma\vdash\langle
        \Lambda 01.C,\mathtt{id},\mathtt{comp}, TT,
        \mathtt{cnst}, \Lambda 01.P,\mathtt{pair},
        \mathtt{proj}_1,\mathtt{proj}_2,
        \Lambda 01.E,\mathtt{lam},\mathtt{eval}
    \rangle\,\mathrm{model}}
\end{mathpar}
\caption{STLC-R declarative typing rules}
\Description{}
\label{stlc-r-typing}
\end{figure}

In figure \ref{stlc-r-typing}, we mostly repeat usual
declarative typing rules of a Hindley-Milner type system
\cite{Clement:ML1986} (TODO: check source) with an
addition of a rule for reinterpretation (discussed
below) and an important change in \texttt{let} case: as
discussed, we don't take local context in account when
checking polymorphic term $s$. Instead, it can only
reference other polymorphic terms. This is safe since
other polymorphic terms are constrained in the same way,
and can be safely used inside $s$.

To check models, we abuse our type system a bit in order
to avoid adding type operators and complicate the type
system: types of arrows $C$, products $P$ and
exponentials $E$ are specified as open monomorphic types
which can use two type variables, $0$ and $1$. This fact
is reflected in the syntax of a model using notation
$\Lambda 01.C$, $\Lambda 01.P$ and $\Lambda 01.E$
respectively. $0$ and $1$ are then filled in by
substitution when we need to type model operations
(notation $T[S_1, S_2]$ is a shorthand for
$T[0\coloneqq S_1,1\coloneqq S_2]$).

We could generalize model definition to include
something else as objects instead of types, but there is
no way to build meaningful $C$,$P$ and $E$ with anything
other than types in base system.

\begin{definition}
    Let $\mathcal{M}=\langle\Lambda 01.C,\mathtt{id},
    \mathtt{comp},TT,\mathtt{cnst},\Lambda 01.P,
    \mathtt{pair},\mathtt{proj}_1,\mathtt{proj}_2,
    \Lambda 01.E,\mathtt{lam},\mathtt{eval}\rangle$.
    Then, interpretation of type in a model is the
    following function:
    \begin{align}{}
    \llbracket T\rrbracket^\mathcal{M}_\sigma &=
        C[TT,T^\mathcal{M}_\sigma] \\
    \alpha^\mathcal{M}_\sigma &= \sigma\alpha \\
    \top^\mathcal{M}_\sigma &= TT \\
    (S\times T)^\mathcal{M}_\sigma &=
        P[S^\mathcal{M}_\sigma,T^\mathcal{M}_\sigma] \\
    (S\to T)^\mathcal{M}_\sigma &=
        E[S^\mathcal{M}_\sigma,T^\mathcal{M}_\sigma].
    \end{align}
\end{definition}

Having defined typing rules, let's pin down computation
rules in order to prove some useful properties. General
idea is the following: before inlining a \texttt{let}
binding, we first have to make sure that the polymorphic
term is closed and is in normal form; in all other
cases, standard reductions of STLC apply.

\begin{figure}
\centering
\begin{align*}
    u,v&\Coloneqq n\;|\;()\;|\;(u,v)\;|\;\lambda x.v\\
    n&\Coloneqq x\;|\;\pi_i\,n\;|\;n\,v
\end{align*}
\caption{Values and neutral terms of STLC-R}
\Description{}
\label{stlc-r-val}
\end{figure}

\begin{figure}
\centering
\begin{mathpar}
    \inferrule{s\to s'}{(s,t)\to(s',t)}
    \and
    \inferrule{t\to t'}{(s,t)\to(s,t')}
    \and
    \inferrule{t\to t'}{\pi_i t\to\pi_i t'}
    \and
    \inferrule{}{\pi_i(t_1, t_2)\to t_i}
    \\
    \inferrule{t\to t'}{\lambda x.t\to\lambda x.t'}
    \and
    \inferrule{s\to s'}{s\,t\to s'\,t}
    \and
    \inferrule{t\to t'}{s\,t\to s\,t'}
    \and
    \inferrule{}{(\lambda x.t)s\to t[x\coloneqq s]}
    \\
    \inferrule{s\to s'}{
        \mathtt{let}\,x=s\,\mathtt{in}\,t\to
        \mathtt{let}\,x=s'\,\mathtt{in}\,t
    }
    \and
    \inferrule{t\to t'}{
        \mathtt{let}\,x=s\,\mathtt{in}\,t\to
        \mathtt{let}\,x=s\,\mathtt{in}\,t'
    }
    \and
    \inferrule{\mathrm{fv}(v)=\varnothing}{
        \mathtt{let}\,x=v\,\mathtt{in}\,t\to
            t\llbracket x\coloneqq v\rrbracket
    }
\end{mathpar}
\caption{STLC-R computation rules}
\Description{}
\label{stlc-r-comp}
\end{figure}

\begin{definition}
    By $t\llbracket x\coloneqq v\rrbracket$ we mean an
    operation which works as a regular substitution of
    $v$ into $t$ except for the case of
    reinterpretation:
    \begin{align}{}
        x^\mathcal{M}\llbracket
            x\coloneqq v\rrbracket &= \llbracket
                v \rrbracket^\mathcal{M}/\varepsilon
        \\
        \llbracket()\rrbracket^\mathcal{M}/\gamma &=
            \mathtt{cnst}
        \\
        \llbracket(s,t)\rrbracket^\mathcal{M}/\gamma &=
            \mathtt{pair}
                \,\llbracket s\rrbracket^\mathcal{M}/
                    \gamma
                \,\llbracket t\rrbracket^\mathcal{M}/
                    \gamma
        \\
        \llbracket\pi_i\,t\rrbracket^\mathcal{M}/\gamma
            &=\mathtt{comp}\,\mathtt{proj}_i\,\llbracket
                t\rrbracket^\mathcal{M}/\gamma
        \\
        \llbracket x\rrbracket^\mathcal{M}/\gamma &=
            \gamma\,x
        \\
        \llbracket\lambda x.t\rrbracket^\mathcal{M}/
            \gamma &= \mathtt{lam}\,\llbracket t
            \rrbracket^\mathcal{M}/\left[
                \vec{y}\coloneqq\mathtt{comp}\,
                    \mathtt{proj}_1(\gamma\,\vec{y}),\;
                x\coloneqq\mathtt{proj}_2
            \right]
        \\
        \llbracket t\,s\rrbracket^\mathcal{M}/\gamma &=
            \mathtt{comp}\,\mathtt{eval}\,\left(
                \mathtt{pair}
                \,\llbracket t\rrbracket^\mathcal{M}/
                    \gamma
                \,\llbracket s\rrbracket^\mathcal{M}/
                    \gamma
            \right)
    \end{align}
    where model operations \texttt{cnst},\,\texttt{pair}
    and others are taken from $\mathcal{M}$ and
    $\varepsilon$ is an empty substitution.
\end{definition}

\subsection{Metatheoretical properties of STLC-R}
\label{STLC-R-properties}

Here we only study some basic properties of a system to
make sure that everything works as expected:
computations don't diverge; types are preserved by all
computations; all of the types can be inferred if you
wish to do so; everything terminates and, modulo
reduction of \texttt{let}-bindings, this is the same old
STLC we started with. Proofs are mostly following
standard techniques, since rules for reinterpretation
are structured to directly plug into proofs of
properties of STLC.

\begin{theorem}[Church-Rosser theorem for STLC-R]
    Reflexive transitive closure of $(\to)$ has
    Church-Rosser property.
\end{theorem}
\begin{proof}
    Classic proof via parallel reduction relation
    \cite{CRproof}. Parallel reduction is defined as:
    \begin{mathpar}
        \inferrule{}{t\twoheadrightarrow t}
        \and
        \inferrule{
            s\twoheadrightarrow s'\\
            t\twoheadrightarrow t'
        }{(s,t)\twoheadrightarrow(s',t')}
        \and
        \inferrule{t\twoheadrightarrow t'}
            {\pi_i\,t\twoheadrightarrow\pi_i\,t'}
        \and
        \inferrule{
            t_1\twoheadrightarrow t_1'\\
            t_2\twoheadrightarrow t_2'
        }
        {\pi_i(t_1,t_2)\twoheadrightarrow t_i'}
        \\
        \inferrule{t\twoheadrightarrow t'}
            {\lambda x.t\twoheadrightarrow\lambda x.t'}
        \and
        \inferrule{
            s\twoheadrightarrow s'\\
            t\twoheadrightarrow t'
        }
        {s\,t\twoheadrightarrow s'\,t'}
        \and
        \inferrule{
            s\twoheadrightarrow s'\\
            t\twoheadrightarrow t'
        }
        {(\lambda x.s)\,t\twoheadrightarrow
            s'[x\coloneqq t']}
        \\
        \inferrule{
            s\twoheadrightarrow s'\\
            t\twoheadrightarrow t'
        }{
            \mathrm{let}\,x=s\,\mathrm{in}\,t
                \twoheadrightarrow
            \mathrm{let}\,x=s'\,\mathrm{in}\,t'
        }
        \and
        \inferrule{t\twoheadrightarrow t'}{
            \mathrm{let}\,x=v\,\mathrm{in}\,t
                \twoheadrightarrow
            t'\llbracket x\coloneqq v\rrbracket
        }
    \end{mathpar}
\end{proof}

\begin{theorem}[Type preservation for STLC-R]
    The following is true:
    \begin{enumerate}
        \item If $\mathcal{M}\,\mathrm{model}$,
            $\gamma\,x:C[\Gamma^\mathcal{M},
            \Gamma(x)^\mathcal{M}]$ for all $x\in\Gamma$
            and $\Gamma\vdash t:T$, then $\llbracket t
            \rrbracket^\mathcal{M}/\gamma:
            C[\Gamma^\mathcal{M},T^\mathcal{M}]$.
        \item If $\Gamma\vdash s:S$ and
            $\Gamma,x:S,\Delta\vdash t:T$, then
            $\Gamma,\Delta\vdash t[x\coloneqq s]:T$.
        \item If $\Gamma\vdash v:\sigma$ and
            $\Gamma,x:\sigma,\Delta\vdash t:T$, then
            $\Gamma,\Delta\vdash t\llbracket
            x\coloneqq v\rrbracket:T$.
        \item If $\Gamma\vdash t:T$ and $t\to t'$, then
            $\Gamma\vdash t':T$.
    \end{enumerate}
\end{theorem}
\begin{proof}
    By structural induction on terms and $(\to)$.
\end{proof}

\begin{theorem}[Strong normalization for STLC-R]
    Closed terms typable in STLC-R are strongly
    normalizing with respect to $(\to)$.
\end{theorem}
\begin{proof}
    By Tait's method of logical relations. Let $SN$ be a
    set of strongly normalizing terms. Let there be the
    following relation $\mathcal{R}\subseteq
    \mathsf{PolyType}\times\mathsf{Term}$:
    \begin{itemize}
        \item $t\in\mathcal{R}_\alpha,\mathcal{R}_\top$
            iff $t\in SN$;
        \item $t\in\mathcal{R}_{T_1\times T_2}$ iff
            $\pi_i\,t\in\mathcal{R}_{T_i}$ for $i=1,2$;
        \item $t\in\mathcal{R}_{S\to T}$ iff for all
            $s\in\mathcal{R}_S$, $t\,s\in\mathcal{R}_T$;
        \item $t\in\mathcal{R}_{\forall\vec{\alpha}.T}$
            iff $t\in SN_\mathtt{let}$ and for all
            $\mathcal{U}\in\mathcal{R}_{univ}$ and
            $\vec{S}$, $nf_\mathtt{let}(t)\,@\,
            \mathcal{U}\in\mathcal{R}_{\mathcal{U}
            \llbracket T\rrbracket[\vec{\alpha}\coloneqq
            \vec{S}]}$.
    \end{itemize}
    $\mathcal{U}\in\mathcal{R}_{univ}$ is defined
    similarly by requiring that all terms inside
    $\mathcal{U}$ belong to respective $\mathcal{R}$-s.

    Then we can prove our main lemma by mutual induction
    on typing rules:
    \begin{itemize}
        \item If $\Gamma\vdash t:T$ and $\gamma$ is a
            substitution s.t. $\gamma\,x\in
            \mathcal{R}_{\Gamma(x)}$ for all $x\in
            \Gamma$ then $\gamma\,t\in\mathcal{R}_T$.
        \item If $\forall\Gamma\vdash t:T$,
            $\vec{\alpha}=fv(T)$ and $\gamma$ is a
            substitution s.t. $\gamma\,x\in
            \mathcal{R}_{\Gamma(x)}$ for all $x\in
            \Gamma$ then $\gamma\,t\in
            \mathcal{R}_{\forall\vec{\alpha}.T}$.
    \end{itemize}
    Noticing $\mathcal{R}_T\subseteq SN$ completes the
    proof.
\end{proof}

\begin{theorem}[Conservativity over STLC]
    Normal forms of closed terms typable in STLC-R are
    typable in STLC.
\end{theorem}
\begin{proof}
    Thanks to type preservation, normal form $nf(t)$ is
    typable in STLC-R. Normal form cannot contain
    let-bindings by definition of $(\to)$, therefore
    each subterm is $\forall$-closed. Therefore $nf(t)$
    has no reinterpretation subterms.
\end{proof}

\begin{theorem}
    Global type inference for STLC-R is decidable.
\end{theorem}
\begin{proof}
    It is enough to extend any unification-based type
    inference algorithm for Hindley-Milner type system
    to support universes and reinterpretation.
    \begin{itemize}
        \item Since universes are provided explicitly,
            it is enough to \textit{check} them.
        \item Since polymorphic functions do not depend
            on local context, their type is fully
            resolved when it's added to context.
        \item Once universe is checked in
            $x\,@\,\mathcal{U}$, we can directly compute
            $\mathcal{U}\llbracket T\rrbracket$ (since
            $\mathcal{U}$ is known and $T$ is fully
            resolved) and instantiate it with fresh
            metavariables in $\vec{S}$.
    \end{itemize}
\end{proof}

\subsection{Useful reinterpretations inside STLC-R}
\label{STLC-R-applications}

The idea to interpret STLC terms in cartesian closed
categories is really not new and is already thoroughly
explored by Conal Elliott in ``Compiling to Categories''
\cite{Elliott:Compcat2017}; our addition here is that we
internalized reinterpretation inside the calculus
without working through compiler plugin. So all of the
main applications from the cited paper are still
applicable here whenever they are expressible as
universes in STLC-R. As STLC-R is not the main focus of
the paper, we will not be reiterating the examples here,
but we would like to add just one more example.

Following Moggi \cite{MOGGI199155}, we can use monads to
model side effects in a language. What is new here is
that instead of parameterizing a function by arbitrary
monad, we can simply \textit{reinterpret} pure terms as
if they were run in monadic context. Therefore, our
universe is a Kleisli category; however, it is only
cartesian, but in general not cartesian \textit{closed}.
The problem is the following operation (here $F$ is a
monad in question, $E$ is whatever exponential object we
were to use):
\[
    \mathtt{lam}:(\alpha\times\beta\to F[\gamma])\to
        (\alpha\to F[E[\gamma,\beta]])
\]
In other words, $F$ has to \textit{distribute} over an
arrow. Well, at least arrow type distributes over
itself:
\begin{align*}
    \mathcal{U}_E &= \langle\\
        &\Lambda01.\,0\to E\to1,\\
        &\lambda x e.x,\\
        &\lambda gfxe.g(fxe)e,\\
        &\top,\\
        &\lambda e.(),\\
        &\Lambda01.\,0\times1,\\
        &\lambda fgxe.(fxe,gxe),\\
        &\lambda pe.\pi_1\,p,\\
        &\lambda pe.\pi_2\,p,\\
        &\Lambda01.\,0\to1,\\
        &\lambda fxey.f(x,y)e,\\
        &\lambda pe.(\pi_1\,p)\,(\pi_2\,p)\\
    \rangle
\end{align*}

Interestingly enough, if, using Curry-Howard
correspondence, we take closed STLC term $t$ to
represent proof of an intuitionistic proposition,
$t\,@\,\mathcal{U}_E$ corresponds to hypothesis
elimination in intuitionistic logic.

We \textit{could} write down a universe for arbitrary
monad if we had capability to interpret terms which do
not use abstractions in arbitrary cartesian category,
which is actually achievable with
\cite{Elliott:Compcat2017}. We leave study of such
``partial'' interpretation for future work.
